% ============================
%   2. Background
% ============================
\section{Theoretical background}
\label{sec:background}

\subsection{Network Intrusion Detection Systems}

A Network Intrusion Detection System inspects network traffic --- either
passively, by tapping into a SPAN port or virtual interface, or actively,
as part of a forwarding plane (NIPS) --- and emits an alert when the
observed packets match a model of malicious activity~\cite{garcia2009anomaly,
buczak2016survey}. NIDS taxonomies typically split detection strategies
along two orthogonal axes:
\begin{description}
    \item[Detection paradigm.] \textit{Misuse} (a.k.a.\
        \textit{signature-based}) detection encodes known attack
        patterns; \textit{anomaly-based} detection encodes normal
        behaviour and reports deviations.
    \item[Granularity.] \textit{Packet-level} detection inspects each
        frame independently; \textit{flow-level} detection aggregates
        packets into bidirectional 5-tuple flows before scoring.
\end{description}

\nidsname{} operates packet-level for the signature engine (rules can
match any combination of L3/L4 fields, payload regular expressions, or
content substrings) and per-packet for the current ML detector, with
flow-level features (inter-arrival time, src-rate) maintained as
streaming state. This trade-off favours real-time responsiveness over
the larger context windows used by flow-based systems such as
Zeek~\cite{paxson1999bro}.

\subsection{Signature-based detection}

The dominant open-source signature systems use a domain-specific
language for rules~\cite{roesch1999snort,suricata}. A typical Snort rule
is structured as a header (action, protocol, addresses, ports) followed
by a comma-separated set of options that constrain payload content,
session state, flow direction and so on. \nidsname{} adopts a simpler,
JSON-based match specification (Listing~\ref{lst:rule}) which is
expressive enough for the categories of attack the project targets
(scans, brute-force, command injection, reverse shell, simple flood)
without requiring users to learn a third-party DSL.

\begin{lstlisting}[language=,caption={Example signature rule, stored both
in YAML and in the database},label={lst:rule}]
- name: "SSH brute force"
  category: "BruteForce"
  severity: "high"
  match:
    protocol: tcp
    dst_port: 22
    flags: "S"
    threshold:
      count: 10
      window_sec: 30
      group_by: src_ip
  action: alert
\end{lstlisting}

\subsection{Machine learning for intrusion detection}

The two algorithms exposed by \nidsname{} were chosen for their
complementary semantics rather than for state-of-the-art performance:
\begin{description}
    \item[Random Forest~\cite{breiman2001random}] is a supervised
        ensemble of decision trees trained on labelled
        $(\text{features}, y)$ pairs. It produces calibrated probability
        estimates, exposes per-feature importance (used to populate the
        dashboard bar chart), and is robust to noisy or irrelevant
        features. This makes it a natural fit for the ``second opinion''
        role described in the introduction.
    \item[Isolation Forest~\cite{liu2008isolation}] is an unsupervised
        anomaly detector that recursively partitions the feature space
        with random splits; points that require few splits to be
        isolated are scored as anomalies. The algorithm does not require
        labelled attack data, which matters in practice because attack
        traffic is far rarer than benign traffic and is hard to label
        reliably~\cite{sommer2010outside}.
\end{description}

The feature vector currently used by both algorithms is
$[\,src\_port,\ dst\_port,\ length,\ proto\_id,\ flag\_mask,\ inter\_arrival\_ms\,]$
--- six numeric features computed per packet, with the inter-arrival
time maintained as streaming state keyed by source IP. This intentionally
small vector keeps inference cheap (sub-millisecond per packet on a
laptop) and avoids the dataset-leakage pitfalls described by Sommer and
Paxson~\cite{sommer2010outside}, where features that encode the
collection methodology rather than the underlying behaviour cause models
to over-fit.

\subsection{Hybrid detection}

When both engines run concurrently, three alert categories are emitted:
\begin{itemize}
    \item \textbf{signature-only} --- the rule engine matched, the ML
        model disagreed (label = benign).
    \item \textbf{ml-only} --- no rule matched, the ML model scored the
        packet above a configurable threshold (currently $0.8$).
    \item \textbf{hybrid} --- both engines agreed. Hybrid alerts carry
        the highest confidence and are surfaced first in the dashboard
        triage queue.
\end{itemize}
This explicit labelling makes the analyst aware of \textit{why} an alert
fired and avoids the black-box critique levelled at pure ML
IDS~\cite{sommer2010outside}.

\subsection{Forensics in a SIEM}

The Security Information and Event Management acronym covers the broader
class of systems that aggregate, correlate and present security
events~\cite{aguilar2017siem}. In a production SIEM, forensics is
typically delegated to an external packet capture appliance. \nidsname{}
takes the opposite approach: every alert triggers a write of the last
$N = 50$ packets seen by the dispatcher (kept in a circular buffer of
size $200$) to a per-alert JSON Lines file. The file is associated with
the alert via a foreign-key relationship in the SQLite database and is
served back to the UI through an authenticated FastAPI endpoint. This
brings the most common triage workflow --- ``show me what happened just
before this alert fired'' --- into the application itself.
